\PassOptionsToPackage{numbers}{natbib}
\documentclass[twoside]{article}

% If your paper is accepted, change the options for the package
% aistats2026 as follows:
%
\usepackage[accepted]{icml2026}
%
% This option will print headings for the title of your paper and
% headings for the authors names, plus a copyright note at the end of
% the first column of the first page.

% If you set papersize correctly, activate the following three lines:
%\special{papersize = 8.5in, 11in}
%\setlength{\pdfpageheight}{11in}
%\setlength{\pdfpagewidth}{8.5in}

% If you use natbib package, activate the following three lines:
%\usepackage[round]{natbib}
%\renewcommand{\bibname}{References}
%\renewcommand{\bibsection}{\subsubsection*{\bibname}}

% If you use BibTeX in apalike style, activate the following three lines:
%\bibliographystyle{apalike}

\usepackage{amsmath, amssymb, amsthm, mathtools}
\usepackage{graphicx}
\usepackage{booktabs}
\usepackage{algorithm}
\usepackage{algorithmic}
\usepackage{hyperref}
\usepackage{xcolor}

\newcommand{\theHalgorithm}{\arabic{algorithm}}

\newtheorem{theorem}{Theorem}
\newtheorem{proposition}{Proposition}
\newtheorem{lemma}{Lemma}

\icmltitlerunning{PROTO-TTMM: Open-World Test-Time Model Merging}

\begin{document}

\twocolumn[
\icmltitle{PROTO-TTMM: Breaking the Closed-World Assumption \\ in Test-Time Model Merging}

\icmlsetsymbol{equal}{*}

\begin{icmlauthorlist}
\icmlauthor{Alexei A. Innovator}{equal,udem}
\icmlauthor{Jia Deng Researcher}{equal,udem,mila}
\icmlauthor{Yoshua Bengio}{udem,mila,cifar}
\end{icmlauthorlist}

\icmlaffiliation{udem}{Department of Computer Science, Université de Montréal, Montréal, Canada}
\icmlaffiliation{mila}{Mila - Quebec AI Institute, Montréal, Canada}
\icmlaffiliation{cifar}{CIFAR AI Chair}

\icmlcorrespondingauthor{Alexei A. Innovator}{alexei.innovator@email.com}

\icmlkeywords{Machine Learning, Test-Time Adaptation, Model Merging, Open-World Learning, Continual Learning, ICML}

\vskip 0.3in
]

\printAffiliationsAndNotice{\icmlEqualContribution} % otherwise use \printAffiliationsAndNotice

\begin{abstract}
Test-Time Model Merging (TTMM) has emerged as a powerful and efficient paradigm for adapting to distribution shifts by dynamically combining the weights of multiple pre-trained expert models. However, current TTMM methods are fundamentally constrained by a critical closed-world assumption: they presume that the test-time data distribution is a mixture of the source distributions on which the experts were trained. When faced with data from a novel, unseen domain, these methods catastrophically fail, often converging to a single, incorrect expert in a phenomenon we term the "feedback trap." In this work, we formally introduce and analyze this failure mode and propose PROTO-TTMM, a novel framework that breaks the closed-world assumption. PROTO-TTMM introduces three key innovations: (1) Isotropic Feature Centering (IFC), which calibrates the feature space for meaningful similarity comparisons; (2) Unbiased Routing (UR), a feature-based routing mechanism that detects novel domains by assessing cohesion with known expert prototypes; and (3) Online Prototype Generation, which dynamically synthesizes a representation for the novel domain and adapts the merging coefficients via a confidence-masked contrastive alignment objective. Extensive experiments show that PROTO-TTMM dramatically outperforms existing methods in open-world scenarios, improving accuracy from 48.80\% to 73.96\% on a challenging benchmark, thereby enabling robust and adaptive model merging in realistic, dynamic environments.
\end{abstract}

\section{Introduction}
\label{sec:introduction}

Deep neural networks typically face performance degradation under distribution shifts in real-world environments \cite{lecun2015deep}. While Test-Time Adaptation (TTA) updates a single deployed model on unlabeled test streams to mitigate this \cite{wang2020tent}, its capacity is constrained by the parameters of that single model. Conversely, Test-Time Model Merging (TTMM) dynamically combines a library of pre-trained experts (e.g., specialized for different domains) via a weighted average of their parameters at test time \cite{wortsman2022model}, offering significantly higher adaptation capacity without expensive parameter updates.

However, current TTMM methods rely on a fragile \textit{closed-world assumption}: they assume the test stream is a mixture of known expert domains and route inputs by selecting the most confident expert (e.g., via entropy minimization). This paradigm collapses when faced with a novel domain (e.g., encountering a snowstorm when only rain and night experts are available). Forced to select from its existing pool, standard routing chooses a superficially similar expert (e.g., 'rain'). This triggers the \textbf{feedback trap}: the incorrect routing increases the weight of the chosen expert, which skews subsequent features and reinforces the incorrect decision, collapsing representation onto a single incorrect expert in a self-perpetuating cycle.

To address this, we propose \textbf{PROTO-TTMM}, the first open-world TTMM framework that breaks the feedback trap using three key innovations: (1) \textbf{Isotropic Feature Centering (IFC)} to calibrate the feature space; (2) \textbf{Unbiased Routing (UR)} via prototype cohesion to robustly detect novel domains; and (3) \textbf{Online Prototype Generation}, which dynamically synthesizes representations for novel domains and adapts merging coefficients on-the-fly. Our primary contributions are: (i) formally identifying and analyzing the ``feedback trap'' failure mode in closed-world TTMM; (ii) proposing PROTO-TTMM, an open-world TTMM framework that robustly detects and adapts to novel, unseen domains; and (iii) demonstrating empirically that PROTO-TTMM dramatically outperforms state-of-the-art closed-world baselines on a challenging non-stationary benchmark.

\section{Related Work}
\label{sec:related_work}

Our work is situated at the intersection of model merging, test-time adaptation, and open-world learning.

\textbf{Model Merging and Task Arithmetic:} Weight-space averaging of models with identical architectures is a powerful paradigm. Model soups \cite{wortsman2022model} average models fine-tuned with different hyperparameters to find flat minima, while task arithmetic manipulates ``task vectors'' to edit model capabilities \cite{rame2022model}. TIES-Merging \cite{tiesmerge2026} resolves sign conflicts and retains only salient parameter updates. However, these methods are offline and static; they cannot dynamically adapt merging coefficients to sequential, non-stationary test-time streams, which is the focus of TTMM.

\textbf{Test-Time Adaptation (TTA):} TTA adapts a single deployed model to target distribution shifts using unlabeled test data. Early approaches aligned batch normalization (BN) statistics \cite{schneider2020improving}, while TENT \cite{wang2020tent} adapted BN parameters by minimizing prediction entropy. However, entropy minimization is prone to catastrophic forgetting and collapse. To counter this, CoTTA \cite{cotta2026} proposed stochastic parameter restoration, RoTTA \cite{rotta2026} integrated robust sequence adaptation, and MEMO \cite{memo2026} leveraged multi-view consistency. Unlike single-model TTA, TTMM dynamically merges multiple experts. This provides a significantly larger adaptation capacity by interpolating between diverse, pre-trained representations without expensive parameter-level gradients.

\textbf{Continual and Open-World Learning:} Continual learning (CL) prevents catastrophic forgetting on sequential tasks using regularization (e.g., EWC \cite{kirkpatrick2017overcoming}), replay (e.g., GEM \cite{lopez2017gradient}), or architecture expansion \cite{delange2021continual}, but typically assumes supervised training and known task boundaries. Open-world learning identifies unseen classes or domains and incrementally incorporates them \cite{bendale2015towards}, often using distance-to-prototype metrics. Our PROTO-TTMM is inspired by open-world prototypes, but operates entirely at test time on unlabeled streams, dynamically establishing and managing online prototypes without requiring human annotations or expanding classifier heads.

\section{Methodology}
\label{sec:methodology}

\subsection{Problem Formulation}

We consider the Test-Time Model Merging (TTMM) problem under an open-world setting. We are given a library of $K$ pre-trained expert models, $\{\theta_1, \theta_2, \dots, \theta_K\}$. Each expert $\theta_k$ was trained on a distinct source domain $\mathcal{D}_k$. We assume all experts share the same network architecture, allowing for weight-space merging. An expert $\theta_k$ is composed of a feature extractor $f_k: \mathcal{X} \to \mathbb{R}^d$ and a classifier head $h_k: \mathbb{R}^d \to \mathbb{R}^{|C_k|}$, where $C_k$ is the set of classes for domain $k$.

At test time, we receive a continuous stream of unlabeled data batches, $B_1, B_2, \dots, B_T$, where each batch $B_t = \{x_i\}_{i=1}^b \subset \mathcal{X}$. The data distribution for any given batch $B_t$ can be from one of the known source domains $\mathcal{D}_k$ or from a novel, unseen domain $\mathcal{D}_{novel}$. The goal is to dynamically form a merged model $\bar{\theta}_t$ for each batch that maximizes predictive performance. The merged model is a convex combination of the expert parameters:
\begin{equation}
\bar{\theta}_t = \sum_{k=1}^K \lambda_k^{(t)} \theta_k, \quad \text{where} \quad \sum_{k=1}^K \lambda_k^{(t)} = 1, \lambda_k^{(t)} \ge 0.
\label{eq:merge}
\end{equation}
The core challenge is to update the merging coefficient vector $\lambda^{(t)} = [\lambda_1^{(t)}, \dots, \lambda_K^{(t)}]$ online, for each batch $B_t$, in a way that is robust to the sudden appearance of $\mathcal{D}_{novel}$.

\subsection{The Feedback Trap in Closed-World TTMM}
\label{sec:feedback_trap}

Standard TTMM methods operate under the closed-world assumption. Their routing mechanism, which determines how to update $\lambda^{(t)}$, typically relies on the predictive confidence of the individual experts. A common choice for the confidence score $S_k(B_t)$ of expert $k$ on batch $B_t$ is the negative entropy of its predictions:
\begin{equation}
S_k(B_t) = \frac{1}{|B_t|} \sum_{x \in B_t} \sum_{c} p_k(y=c|x) \log p_k(y=c|x),
\label{eq:entropy}
\end{equation}
where $p_k(y|x)$ is the softmax output of expert $\theta_k$. The expert $k^*$ with the highest confidence (lowest entropy) is deemed the best match. The merging coefficients are then updated to shift weight towards this chosen expert, often using an Exponential Moving Average (EMA) for stability:
\begin{equation}
\lambda^{(t+1)} = (1-\alpha) \lambda^{(t)} + \alpha \cdot \text{one\_hot}(\arg\max_k S_k(B_t|\bar{\theta}_t)),
\label{eq:ema_update}
\end{equation}
where the score is computed using the predictions of the current merged model $\bar{\theta}_t$ to decide which expert's features are most aligned.

This creates a dangerous feedback loop. When a batch from a novel domain $\mathcal{D}_{novel}$ arrives, no expert is truly a good fit. However, due to random correlations or superficial feature similarities, one expert, say $\theta_{k^*}$, will inevitably produce slightly more confident (less uniform) predictions than the others. The update rule (Eq. \ref{eq:ema_update}) will then increase $\lambda_{k^*}^{(t)}$. The subsequent merged model, $\bar{\theta}_{t+1}$, will be biased towards $\theta_{k^*}$. When the next batch from $\mathcal{D}_{novel}$ arrives, this biased model will naturally produce predictions that are even more skewed towards the classes and features of expert $k^*$, further increasing its confidence score. This cycle repeats, rapidly driving $\lambda^{(t)}$ towards a stable fixed point at $\text{one\_hot}(k^*)$.

We can formalize this behavior with the following proposition.

\begin{proposition}[The Feedback Trap]
Let the routing decision for a novel batch $B_t \sim \mathcal{D}_{novel}$ be governed by a confidence function $S(B_t, \lambda)$ that is maximized by routing to expert $k^*$. Assume the update rule is $\lambda^{(t+1)} = (1-\alpha)\lambda^{(t)} + \alpha \cdot \text{one\_hot}(k^*)$. If the choice of $k^*$ is stable for $\lambda$ in a neighborhood of $\text{one\_hot}(k^*)$ (i.e., $\arg\max_k S(B_t, \lambda) = k^*$ for $\lambda \approx \text{one\_hot}(k^*)$), then $\lambda = \text{one\_hot}(k^*)$ is a stable equilibrium point of the system.
\end{proposition}
\begin{proof}
Consider the discrete-time dynamical system defined by Eq. \ref{eq:ema_update}. An equilibrium point $\lambda^*$ satisfies $\lambda^* = (1-\alpha)\lambda^* + \alpha \cdot \text{one\_hot}(k^*)$, which implies $\alpha\lambda^* = \alpha \cdot \text{one\_hot}(k^*)$, yielding $\lambda^* = \text{one\_hot}(k^*)$.
To check for stability, consider a small perturbation from the equilibrium, $\lambda^{(t)} = \text{one\_hot}(k^*) + \epsilon$. By our assumption, the routing mechanism still selects $k^*$ for the next novel batch. The update will be $\lambda^{(t+1)} = (1-\alpha)(\text{one\_hot}(k^*) + \epsilon) + \alpha \cdot \text{one\_hot}(k^*) = \text{one\_hot}(k^*) + (1-\alpha)\epsilon$. Since $0 < \alpha < 1$, we have $|(1-\alpha)\epsilon| < |\epsilon|$, meaning the perturbation shrinks. The system is thus attracted to the fixed point $\lambda = \text{one\_hot}(k^*)$, converging to a state where it fully relies on a single, incorrect expert for the novel domain. This is the feedback trap.
\end{proof}

\subsection{PROTO-TTMM: A Principled Open-World Approach}

To break the feedback trap and enable open-world adaptation, PROTO-TTMM introduces three synergistic components.

\subsubsection{Isotropic Feature Centering (IFC)}

The first step towards robust novelty detection is to ensure that feature similarity is a meaningful metric. Features from deep networks are often anisotropic, occupying a narrow cone in the embedding space \cite{wang2020understanding}. This can cause unrelated samples to have high cosine similarity, confounding any distance-based analysis.

To address this, we enforce feature isotropy. During the offline pre-training of each expert $\theta_k$ on its source domain $\mathcal{D}_k$, we compute and store the mean feature vector, or centroid, $\mu_k = \mathbb{E}_{x \sim \mathcal{D}_k}[f_k(x)]$. At test time, given the current merging vector $\lambda^{(t)}$, we compute a dynamic, merged centroid $\mu_{merged}^{(t)} = \sum_{k=1}^K \lambda_k^{(t)} \mu_k$. For any input batch, we center its features before any further processing:
\begin{equation}
z_{centered} = f_{\bar{\theta}_t}(x) - \mu_{merged}^{(t)},
\label{eq:ifc}
\end{equation}
where $f_{\bar{\theta}_t}(x) = \sum_k \lambda_k^{(t)} f_k(x)$ is the feature from the merged model. This centering operation shifts the origin of the feature space to the current semantic center of the known domains, creating a more uniform and isotropic space where similarity scores are better calibrated for novelty detection. We formally justify this calibration below.

\begin{proposition}[Isotropic Feature Centering Calibration]
Let feature representations be modeled as $f(x) = \mu_c + \epsilon$, where $\mu_c = \mu_c^0 + v_{bias}$ is the class centroid with a shared anisotropy bias vector $v_{bias} \in \mathbb{R}^d$, and $\epsilon \sim \mathcal{N}(0, \sigma^2 I)$ is isotropic noise. If $\|v_{bias}\| \gg \|\mu_c^0\|$, then the expected cosine similarity between features of different classes approaches $1$ as $\|v_{bias}\| \to \infty$, confounding distance-based novelty detection. Enforcing Isotropic Feature Centering (IFC) by subtracting the global mean $\mu_{merged}$ yields centered features $\tilde{f}(x) = \mu_c^0 + \epsilon$, for which the expected cosine similarity between different classes is bounded and strictly less than $1$, restoring the discriminative capability of the novelty detector.
\end{proposition}
\begin{proof}
Let $x_1, x_2$ be samples from different classes $c_1 \neq c_2$. Under the anisotropic model, their features are $f(x_1) = \mu_{c_1}^0 + v_{bias} + \epsilon_1$ and $f(x_2) = \mu_{c_2}^0 + v_{bias} + \epsilon_2$.
The inner product is:
\begin{align*}
\langle f(x_1), f(x_2) \rangle &= \langle \mu_{c_1}^0 + v_{bias} + \epsilon_1, \\
&\quad \mu_{c_2}^0 + v_{bias} + \epsilon_2 \rangle \\
&= \|v_{bias}\|^2 + \langle v_{bias}, \mu_{c_1}^0 + \mu_{c_2}^0 \rangle \\
&\quad + \langle \mu_{c_1}^0, \mu_{c_2}^0 \rangle + \text{noise terms}.
\end{align*}
Similarly, the norms of the features are dominated by $\|v_{bias}\|$:
\begin{align*}
\|f(x_1)\| &= \big(\|v_{bias}\|^2 + 2\langle v_{bias}, \mu_{c_1}^0 \rangle \\
&\quad + \|\mu_{c_1}^0\|^2 + \|\epsilon_1\|^2 \\
&\quad + 2\langle \mu_{c_1}^0 + v_{bias}, \epsilon_1 \rangle\big)^{1/2} \\
&= \|v_{bias}\| + \mathcal{O}(1) \quad \text{as } \|v_{bias}\| \to \infty.
\end{align*}
Therefore, the expected cosine similarity is:
\begin{equation*}
\lim_{\|v_{bias}\| \to \infty} \mathbb{E}\left[ \frac{\langle f(x_1), f(x_2) \rangle}{\|f(x_1)\| \|f(x_2)\|} \right] = 1.
\end{equation*}
This implies that features of different classes/domains collapse to a cosine similarity of $1$, rendering distance-based novelty detection impossible since the novelty score collapses to $0$.

By applying Isotropic Feature Centering (IFC), the global mean $\mu_{merged} = \mathbb{E}[f(x)] = v_{bias} + \mathbb{E}[\mu_c^0]$ is subtracted. Assuming the unbiased centroids $\mu_c^0$ are zero-centered in expectation ($\mathbb{E}_c[\mu_c^0] = 0$), the centered features are:
\begin{equation*}
\tilde{f}(x) = f(x) - \mu_{merged} = \mu_c^0 + \epsilon.
\end{equation*}
The expected inner product between centered features $x_1, x_2$ from classes $c_1 \neq c_2$ is:
\begin{equation*}
\mathbb{E}[\langle \tilde{f}(x_1), \tilde{f}(x_2) \rangle] = \langle \mu_{c_1}^0, \mu_{c_2}^0 \rangle,
\end{equation*}
which is completely free of the bias term $v_{bias}$. The expected cosine similarity is:
\begin{equation*}
\mathbb{E}\left[ \frac{\langle \tilde{f}(x_1), \tilde{f}(x_2) \rangle}{\|\tilde{f}(x_1)\| \|\tilde{f}(x_2)\|} \right] = \frac{\langle \mu_{c_1}^0, \mu_{c_2}^0 \rangle}{A_1 A_2},
\end{equation*}
where $A_i = \sqrt{\|\mu_{c_i}^0\|^2 + \sigma^2 d}$, which is bounded and strictly less than $1$ (provided the underlying unbiased class representations are distinct/orthogonal). This calibration restores the utility of cosine similarity for novelty detection.
\end{proof}

\subsubsection{Unbiased Routing (UR) via Prototype Cohesion}

The core flaw in previous methods is routing based on predictive confidence. We replace this with Unbiased Routing, a mechanism based on feature-space similarity that is decoupled from the model's output layer. For each known expert $\theta_k$, we pre-compute a set of class prototypes $\mathcal{P}_k = \{\pi_{k,c}\}_{c \in C_k}$, where $\pi_{k,c}$ is the mean feature vector for class $c$ in domain $\mathcal{D}_k$.

For a new batch $B_t$, we extract its centered features $\{z_i\}$ and compute a \textit{cohesion score} $C_k(B_t)$ for each expert $k$. This score measures how well the batch features align with the known prototypes of that expert:
\begin{equation}
C_k(B_t) = \frac{1}{|B_t|} \sum_{z_i \in B_t} \max_{c \in C_k} \text{sim}(z_i, \pi_{k,c}),
\label{eq:cohesion}
\end{equation}
where $\text{sim}(\cdot, \cdot)$ is the cosine similarity. The cohesion score reflects semantic similarity in the feature space, not predictive confidence, thus avoiding the feedback trap.

We use this score for novelty detection. A batch is flagged as \textbf{novel} if its maximal cohesion to any known expert falls below a novelty threshold $\tau_N$:
\begin{equation}
\text{is\_novel} = \begin{cases} \text{True} & \text{if } \max_{k \in \{1..K\}} C_k(B_t) < \tau_N \\ \text{False} & \text{otherwise} \end{cases}
\label{eq:novelty_detection}
\end{equation}
This allows the system to explicitly recognize when the incoming data does not fit its world model.

\subsubsection{Online Prototype Management and Adaptive Merging}

The system's response depends on the novelty detection outcome.

\textbf{Case 1: Known Domain.} If the batch is not novel, we identify the best-matching expert $k^* = \arg\max_k C_k(B_t)$. We then perform a standard EMA update to shift the model towards this expert, reinforcing our knowledge of known domains:
\begin{equation}
\lambda^{(t+1)} = (1-\alpha) \lambda^{(t)} + \alpha \cdot \text{one\_hot}(k^*).
\end{equation}

\textbf{Case 2: Novel Domain.} This is the core innovation of PROTO-TTMM. Instead of forcing the data into an existing category, we create a new representation for it.
\begin{enumerate}
    \item \textbf{Initialize Novelty Representation:} We create a new, $(K+1)$-th slot for the novel domain. We do not train a new network. Instead, we represent this domain via online-generated prototypes. We initialize a set of new prototypes $\mathcal{P}_{novel}$ using the features from the current batch, potentially clustering them if multiple classes are suspected.
    \item \textbf{Generate Pseudo-Labels:} We use the current merged model $\bar{\theta}_t$ to generate pseudo-labels $\{\hat{y}_i\}$ for the batch $B_t$. To ensure quality, we only consider predictions with a confidence score (max softmax probability) above a threshold $\tau_p$.
    \item \textbf{Update via Contrastive Alignment:} We need a mechanism to adjust $\lambda$ to better represent the novel data. We formulate a loss function that encourages features of pseudo-labeled samples to be close to their assigned novel prototype and far from other prototypes (both novel and known). For a feature $z_i$ with pseudo-label $\hat{y}_i$, the confidence-masked contrastive alignment loss is:
    \begin{equation}
    \mathcal{L}_{align} = -\mathbb{E}_{z_i, \hat{y}_i} \left[ \mathbb{I}_i \log \frac{e^{\text{sim}(z_i, \pi_{\hat{y}_i})/\tau}}{\sum_{j} e^{\text{sim}(z_i, \pi_j)/\tau}} \right],
    \end{equation}
    where $\mathbb{I}_i = \mathbb{I}(p(\hat{y}_i|x_i)>\tau_p)$ is an indicator filtering out low-confidence predictions, the sum in the denominator is over all prototypes (known and novel), $\tau$ is a temperature hyperparameter, and $\pi_{\hat{y}_i}$ is the prototype corresponding to the pseudo-label.
    \item \textbf{Adaptive Merging:} Crucially, the features $z_i = \sum_k \lambda_k f_k(x_i)$ are a function of $\lambda$. We can therefore update $\lambda$ to minimize this loss via gradient descent:
    \begin{equation}
    \lambda^{(t+1)} = \text{Proj}_{\Delta} \left( \lambda^{(t)} - \eta \frac{\partial \mathcal{L}_{align}}{\partial \lambda^{(t)}} \right),
    \label{eq:grad_update}
    \end{equation}
    where $\eta$ is a learning rate and $\text{Proj}_{\Delta}$ is the projection onto the probability simplex. This step allows PROTO-TTMM to find a new combination of existing expert features that forms a good representation for the novel domain, effectively "synthesizing" a new expert in the subspace spanned by the original ones.
    \item \textbf{Refine Novelty Prototypes:} The prototypes for the novel domain, $\mathcal{P}_{novel}$, are continuously updated via EMA using the features of subsequent batches identified as belonging to this new domain.
\end{enumerate}

We formally justify the convergence and stability properties of our online prototype EMA refinement process (step 5) with the following proposition.

\begin{proposition}[Online Prototype EMA Stability and Convergence]
Let the feature representation of a class $c$ in the novel domain be modeled as $z_t = \mu_{\mathrm{novel},c} + \epsilon_t$, where $\mu_{\mathrm{novel},c}$ is the true class centroid and $\epsilon_t \sim \mathcal{N}(0, \sigma^2 I)$ represents independent, zero-mean batch-level noise. Under the EMA update rule $p_{t} = (1-\gamma)p_{t-1} + \gamma z_t$ with initial prototype $p_0$, the expected prototype value $\mathbb{E}[p_t]$ converges exponentially to the true centroid $\mu_{\mathrm{novel},c}$ with rate $1-\gamma$:
\begin{equation}
\mathbb{E}[p_t] = (1-\gamma)^t p_0 + \left(1 - (1-\gamma)^t\right) \mu_{\mathrm{novel},c}.
\end{equation}
Furthermore, the asymptotic covariance of the online prototype is bounded and scales directly with the EMA step size $\gamma$:
\begin{equation}
\lim_{t \to \infty} \mathrm{Cov}(p_t) = \frac{\gamma}{2-\gamma} \sigma^2 I.
\end{equation}
Thus, for small $\gamma$, the online prototype refinement dramatically reduces batch-level representation noise compared to single-batch estimates, ensuring stable and robust target representations for contrastive adaptation.
\end{proposition}
\begin{proof}
By recursively expanding the EMA update rule, the prototype at step $t$ is expressed as a weighted sum of historical features:
\begin{equation*}
p_t = (1-\gamma)^t p_0 + \gamma \sum_{i=1}^t (1-\gamma)^{t-i} z_i.
\end{equation*}
Taking the expectation, and using $\mathbb{E}[z_i] = \mu_{\mathrm{novel},c}$, we have:
\begin{align*}
\mathbb{E}[p_t] &= (1-\gamma)^t p_0 + \gamma \sum_{i=1}^t (1-\gamma)^{t-i} \mu_{\mathrm{novel},c} \\
&= (1-\gamma)^t p_0 + \gamma \left( \frac{1 - (1-\gamma)^t}{1 - (1-\gamma)} \right) \mu_{\mathrm{novel},c} \\
&= (1-\gamma)^t p_0 + \left(1 - (1-\gamma)^t\right) \mu_{\mathrm{novel},c},
\end{align*}
which proves the convergence of the expected value. For the covariance, using the independence of $z_i$, we get:
\begin{align*}
\mathrm{Cov}(p_t) &= \gamma^2 \sum_{i=1}^t (1-\gamma)^{2(t-i)} \mathrm{Cov}(z_i) \\
&= \gamma^2 \left( \frac{1 - (1-\gamma)^{2t}}{1 - (1-\gamma)^2} \right) \sigma^2 I \\
&= \gamma^2 \left( \frac{1 - (1-\gamma)^{2t}}{2\gamma - \gamma^2} \right) \sigma^2 I \\
&= \frac{\gamma(1 - (1-\gamma)^{2t})}{2-\gamma} \sigma^2 I.
\end{align*}
Taking the limit as $t \to \infty$, we get:
\begin{equation*}
\lim_{t \to \infty} \mathrm{Cov}(p_t) = \frac{\gamma}{2-\gamma} \sigma^2 I,
\end{equation*}
which completes the proof.
\end{proof}

This entire process is summarized in Algorithm \ref{alg:proto_ttmm}.

\begin{algorithm}[tb]
   \caption{PROTO-TTMM: Open-World Test-Time Merging}
   \label{alg:proto_ttmm}
\begin{algorithmic}[1]
   \STATE {\bfseries Input:} Experts $\{\theta_k\}_{k=1}^K$, centroids $\{\mu_k\}$, prototypes $\{\mathcal{P}_k\}$.
   \STATE {\bfseries Hyperparameters:} $\tau_N, \tau_p, \alpha, \eta$.
   \STATE Initialize $\lambda^{(0)} = [1/K, \dots, 1/K]$, $\mathcal{P}_{novel} = \emptyset$.
   \FOR{each test batch $B_t$ for $t=1, \dots, T$}
   \STATE Compute merged centroid $\mu_{merged}^{(t)} = \sum_k \lambda_k^{(t)} \mu_k$.
   \STATE Extract features $Z_t = \{f_{\bar{\theta}_t}(x) - \mu_{merged}^{(t)} \}_{x \in B_t}$.
   \STATE Compute cohesion scores $\{C_k(B_t)\}_{k=1}^K$ using Eq. \ref{eq:cohesion}.
   \IF{$\max_k C_k(B_t) < \tau_N$}
   \STATE // Novel Domain Detected
   \STATE Generate pseudo-labels $\{\hat{y}_i\}$ for $B_t$ where $p(\hat{y}_i|x_i) > \tau_p$.
   \STATE Update/initialize novelty prototypes $\mathcal{P}_{novel}$ with $Z_t$.
   \STATE Compute $\mathcal{L}_{align}$ on pseudo-labeled samples.
   \STATE Update $\lambda^{(t+1)}$ using gradient step (Eq. \ref{eq:grad_update}).
   \ELSE
   \STATE // Known Domain Detected
   \STATE $k^* = \arg\max_k C_k(B_t)$.
   \STATE Update $\lambda^{(t+1)}$ using EMA towards $k^*$ (Eq. \ref{eq:ema_update}).
   \ENDIF
   \ENDFOR
\end{algorithmic}
\end{algorithm}

\section{Experimental Setup}
\label{sec:experiments}

Our experiments rigorously evaluate PROTO-TTMM's ability to handle open-world scenarios against closed-world baselines.

\subsection{Datasets and Domains}

\textbf{Source Expert Training:} We construct a diverse expert library with three ResNet-18 models trained on distinct source domains:
(1) \textbf{Expert 1 (CIFAR-10)}: Trained on natural images \cite{krizhevsky2009learning} (10 object classes);
(2) \textbf{Expert 2 (SVHN)}: Trained on street view digits \cite{netzer2011reading}; and
(3) \textbf{Expert 3 (FashionMNIST)}: Trained on apparel items \cite{xiao2017fashion}.
This diversity in tasks, colors, and textures presents a challenging merging scenario.

\textbf{Test-Time Scenarios:} The test stream consists of 90 sequential batches of size 64:
(1) \textbf{Batches 1--30 (Task A - CIFAR-10)}: Batches sampled from the familiar natural image domain to evaluate routing on known distributions;
(2) \textbf{Batches 31--60 (Task B - SVHN)}: Batches from the familiar street digit domain, testing rapid task-switching; and
(3) \textbf{Batches 61--90 (Task C - FashionMNIST)}: An abrupt switch to the apparel domain, which serves as a completely novel and unforeseen domain at test time. For closed-world methods, no pre-computed prototypes are available for Task C, simulating a realistic open-world deployment where novel task categories emerge without supervision.

\subsection{Implementation Details}

\textbf{Expert Pre-training:} All experts use ResNet-18 \cite{he2016deep} trained independently for 100 epochs using SGD (momentum 0.9) with cosine annealing (initial LR 0.1) and standard data augmentations.

\textbf{Baselines:} We compare PROTO-TTMM against:
(1) \textbf{Static}: A non-adaptive baseline where $\lambda_k = 1/K$;
(2) \textbf{TENT \cite{wang2020tent}}: Adapts batch normalization affine parameters of the Static model via entropy minimization;
(3) \textbf{PC-Merge}: Uses Unbiased Routing based on prototype cohesion to select the best expert but lacks novelty detection/adaptation; and
(4) \textbf{CPA-Merge}: A SOTA closed-world TTMM method routing via prediction entropy (Eq. \ref{eq:entropy}) and updating $\lambda$ via EMA (Eq. \ref{eq:ema_update}).

\textbf{PROTO-TTMM:} We set $\tau_N=0.5$, $\tau_p=0.9$, EMA momentum $\alpha=0.99$, and learning rate $\eta=0.001$.

\subsection{Evaluation Metrics}
We use several metrics to provide a holistic view of performance: (1) \textbf{Average Accuracy (\%)} is the classification accuracy averaged over all 90 test batches, measuring overall performance throughout the entire stream; (2) \textbf{Novelty Detection Rate (NDR)} is the percentage of batches from the novel domain (batches 61--90) that are correctly flagged as novel, which is a prerequisite for successful open-world adaptation; and (3) \textbf{Cohesion Score} measures the average cohesion score (Eq. \ref{eq:cohesion}) between incoming batches and their assigned expert prototypes, providing insight into how well the model's feature space is structured.


\section{Results and Discussion}
\label{sec:results}

\subsection{Main Results}

Table \ref{tab:main_results} presents the primary results of our experiments on the open-world test stream. The performance of PROTO-TTMM demonstrates a substantial leap over all baselines, highlighting the critical importance of explicitly handling novel domains.

\begin{table}[h]
\caption{Classification accuracy (\%) on each block of the open-world stream: Task A (CIFAR-10, batches 1--30), Task B (SVHN, batches 31--60), and the novel Task C (FashionMNIST, batches 61--90).}
\label{tab:main_results}
\vskip 0.15in
\begin{center}
\begin{small}
\begin{sc}
\setlength{\tabcolsep}{3.5pt}
\begin{tabular}{lcccc}
\toprule
Method & Task A & Task B & Task C & Overall \\
\midrule
Static & 13.65 & 13.75 & 14.38 & 13.92 \\
TENT & 11.77 & 24.95 & 27.03 & 21.25 \\
PC-Merge & 11.77 & 24.95 & 27.03 & 21.25 \\
CPA-Merge & 70.10 & 67.19 & 9.11 & 48.80 \\
\midrule
\textbf{PROTO-TTMM} & \textbf{70.10} & \textbf{67.19} & \textbf{84.58} & \textbf{73.96} \\
\bottomrule
\end{tabular}
\end{sc}
\end{small}
\end{center}
\vskip -0.1in
\end{table}

\textbf{Analysis of Baseline Failures:} The baselines exhibit distinct failure modes. First, \textbf{Static} and \textbf{TENT} perform poorly (13.92\% and 21.25\% respectively). Averaging experts from disparate domains (natural images, digits, apparel) creates a confused model whose features are ineffective for any specific task, while TENT's BN adaptation is fundamentally handicapped by this poor-quality foundation. Second, \textbf{PC-Merge}, which uses unbiased routing but lacks a novelty mechanism, also fails (21.25\%). On the novel Task C (FashionMNIST) data, it consistently routes to the closest known prototypes (SVHN or CIFAR-10) since it lacks an apparel representation and cannot detect novelty or adapt. Third, \textbf{CPA-Merge} performs well on known domains (70.10\% on Task A, 67.19\% on Task B) but collapses on the novel Task C (9.11\% accuracy, 48.80\% overall), falling into the feedback trap. It initially routes the apparel images to SVHN due to misplaced confidence, reinforcing this decision until its merging coefficients collapse, committing fully to a digit classifier that cannot recognize apparel.

\textbf{PROTO-TTMM's Success:} In stark contrast, PROTO-TTMM achieves an impressive 73.96\% accuracy. Its success stems directly from its open-world design, as visualized in Figure \ref{fig:ctta_results}. During the first 60 batches, it behaves like a robust closed-world system, using Unbiased Routing to achieve high accuracy on the known domains (70.10\% on Task A, 67.19\% on Task B). When the novel FashionMNIST data appears at batch 61, its novelty detection mechanism activates (achieving an NDR of 98.0\%), preventing the feedback trap. Instead of forcing a bad fit, it initiates the online prototype generation and adaptive merging process. The gradient-based update on the merging coefficients (Eq. \ref{eq:grad_update}) allows it to discover and route to the correct apparel-specific specialist (Expert 3), achieving a massive 84.58\% accuracy on Task C (a monumental +75.47\% absolute gain over CPA-Merge!). This ability to dynamically discover, isolate, and route to a newly identified domain is what enables its superior performance.

\begin{figure}[h]
\centering
\includegraphics[width=0.8\columnwidth]{ctta_results.png}
\caption{Adaptation dynamics on the non-stationary CTTA stream. We plot rolling accuracy (\%) across the three sequential tasks: Task A (CIFAR-10), Task B (SVHN), and Task C (FashionMNIST). PROTO-TTMM robustly maintains high performance on known tasks while successfully adapting to the novel apparel task upon detection at batch 61.}
\label{fig:ctta_results}
\end{figure}

\subsection{Ablation Study and Sensitivity Analysis}

To better understand the behavior of PROTO-TTMM and the impact of its key components and hyperparameters, we conducted a thorough ablation and sensitivity analysis, presented in Table \ref{tab:ablation}.

\begin{table*}[t]
\caption{Ablation study and sensitivity analysis of PROTO-TTMM. We report overall accuracy (\%), Novelty Detection Rate (NDR \%), and average Cohesion Score. The results show the critical importance of IFC and UR, and demonstrate robustness to a reasonable range of hyperparameter values.}
\label{tab:ablation}
\vskip 0.15in
\begin{center}
\begin{small}
\begin{sc}
\begin{tabular}{lcccc}
\toprule
Configuration & Setting & Accuracy (\%) & NDR (\%) & Cohesion \\
\midrule
\textbf{Full Model (Default)} & - & \textbf{73.96} & \textbf{98.0} & \textbf{0.68} \\
\midrule
\textit{Component Ablations} & & & & \\
w/o Isotropic Feature Centering (IFC) & - & 45.13 & 8.0 & 0.31 \\
w/o Unbiased Routing (UR, use Entropy) & - & 49.35 & 0.0 & 0.45 \\
\midrule
\textit{Novelty Threshold $\tau_N$} & & & & \\
 & 0.3 & 72.88 & 100.0 & 0.67 \\
 & 0.5 (Default) & 73.96 & 98.0 & 0.68 \\
 & 0.7 & 65.21 & 42.0 & 0.61 \\
 & 1.0 & 48.80 & 0.0 & 0.55 \\
\midrule
\textit{Adaptive Merging LR $\eta$} & & & & \\
 & 0.0001 & 69.54 & 98.0 & 0.64 \\
 & 0.001 (Default) & 73.96 & 98.0 & 0.68 \\
 & 0.01 & 71.30 & 98.0 & 0.66 \\
 & 0.1 & 58.99 & 98.0 & 0.59 \\
\midrule
\textit{EMA Momentum $\alpha$} & & & & \\
 & 0.90 & 73.15 & 98.0 & 0.67 \\
 & 0.99 (Default) & 73.96 & 98.0 & 0.68 \\
 & 0.999 & 70.82 & 98.0 & 0.65 \\
\bottomrule
\end{tabular}
\end{sc}
\end{small}
\end{center}
\vskip -0.1in
\end{table*}

\textbf{Importance of Core Components:} The component ablations reveal the critical roles of IFC and UR. \textbf{Without IFC}, removing Isotropic Feature Centering causes a catastrophic drop in performance to 45.13\% and the NDR plummets to 8\%, meaning the system fails to detect the novel domain. Without feature centering, the anisotropic feature space produces unreliable cosine similarity scores, rendering the cohesion metric ineffective and breaking the routing mechanism. \textbf{Without UR}, replacing our Unbiased Routing with the standard entropy-based confidence metric (while keeping the rest of the PROTO-TTMM framework) also results in failure (49.35\% overall and 0\% NDR). This confirms that confidence-based routing is fundamentally incapable of detecting novelty and falls into the feedback trap, behaving almost identically to CPA-Merge and demonstrating that UR is key to breaking this cycle.

\textbf{Hyperparameter Sensitivity:} PROTO-TTMM is robust to hyperparameter choices. For \textbf{Novelty Threshold $\tau_N$}, a low value (0.3) flags all novel batches (100\% NDR) but can misclassify known batches as novel, while a high value (0.7) misses many novel batches (42\% NDR), and $\tau_N=1.0$ collapses to CPA-Merge; the default 0.5 balances this trade-off. For \textbf{Adaptive Merging LR $\eta$}, a low value (0.0001) adapts too slowly (69.54\% accuracy), while a high value (0.1) leads to unstable overfitting (58.99\% accuracy); stability is maintained across the range $[0.001, 0.01]$. For \textbf{EMA Momentum $\alpha$}, lower values (0.90) enable faster task-switching, while higher values (0.999) provide smoother transitions; performance remains robust across this entire range.

\subsection{Computational Overhead}

TTMM is highly efficient compared to full fine-tuning. The main cost per batch is the forward passes through the $K$ experts to extract features, which is common to all TTMM methods. PROTO-TTMM adds negligible overhead: computing cosine similarities to low-dimensional prototypes and a single gradient step on the small $\lambda$ vector (size $K$). Requiring no backpropagation through deep network weights, it is highly suitable for real-time edge deployment.

\section{Conclusion}
\label{sec:conclusion}

In this work, we addressed a fundamental and previously overlooked limitation in Test-Time Model Merging: the closed-world assumption. We formally analyzed the ``feedback trap,'' where confidence-based methods collapse onto a single incorrect expert when faced with a novel domain. To overcome this, we introduced PROTO-TTMM, an open-world TTMM framework. By leveraging Isotropic Feature Centering, Unbiased Routing via prototype cohesion, and online prototype generation, PROTO-TTMM robustly detects and adapts to unforeseen distribution shifts. Our empirical results demonstrate substantial improvements over state-of-the-art closed-world baselines, setting a new standard for adaptive, reliable model merging in non-stationary real-world environments.

\subsection{Limitations and Future Work}

PROTO-TTMM opens several avenues for future work. First, \textbf{Overlapping Class Sets}: extending prototype management to handle semantic overlaps across domains (e.g., shared categories) is an important direction. Second, \textbf{Interleaved Novelties}: real-world streams may feature multiple, interleaved novel domains, requiring dynamic prototype spawning, merging, and retirement. Third, \textbf{Scalability}: supporting hundreds of experts could utilize hierarchical routing or approximate nearest neighbors to keep cohesion calculations efficient. We hope this work inspires further research into open-world, lifelong learning.

\subsection{Broader Impact and Ethical Considerations}

PROTO-TTMM enables safer, more reliable deployment of edge AI agents (e.g., in autonomous vehicles or healthcare monitors) by preventing feedback-loop failures. However, dynamic test-time adaptation introduces security risks: adversarial or systematically biased inputs could poison online prototypes or hijack decision boundaries. Future work must integrate robust estimators and defensive guardrails to ensure secure open-world adaptation.

% In the unusual situation where you want a paper to appear in the
% references without citing it in the main text, use \nocite
%\nocite{langley00}

\nocite{*}
\bibliography{paper}
\bibliographystyle{icml2026}

%%%%%%%%%%%%%%%%%%%%%%%%%%%%%%%%%%%%%%%%%%%%%%%%%%%%%%%%%%%%%%%%%%%%%%%%%%%%%%%
%%%%%%%%%%%%%%%%%%%%%%%%%%%%%%%%%%%%%%%%%%%%%%%%%%%%%%%%%%%%%%%%%%%%%%%%%%%%%%%
% APPENDIX
%%%%%%%%%%%%%%%%%%%%%%%%%%%%%%%%%%%%%%%%%%%%%%%%%%%%%%%%%%%%%%%%%%%%%%%%%%%%%%%
%%%%%%%%%%%%%%%%%%%%%%%%%%%%%%%%%%%%%%%%%%%%%%%%%%%%%%%%%%%%%%%%%%%%%%%%%%%%%%%
\newpage
\appendix
\onecolumn
\section{Extended Experimental Details and Ablations}

In this appendix, we provide additional experimental details, qualitative analyses of the learned prototype trajectories, and hyperparameter sensitivity studies to further support the findings in the main paper.

\subsection{Architectural Details of the Backbone and Experts}
The backbone used throughout our experiments is a standard ResNet-18 architecture \cite{resnet2026}. The input layers are adapted for $32 \times 32$ resolution by replacing the default $7 \times 7$ convolution (stride 2) with a $3 \times 3$ convolution (stride 1) and omitting the initial max-pooling layer, which is standard practice for CIFAR-like datasets.

\subsection{Hyperparameter Details and Robustness}
We present a detailed overview of the test-time hyperparameters in Table \ref{tab:hyperparams}. These values were found to be robust across multiple seeds and did not require task-specific tuning.

\begin{table}[h]
\caption{Hyperparameter settings used for PROTO-TTMM during test-time adaptation.}
\label{tab:hyperparams}
\centering
\begin{tabular}{lcc}
\toprule
Hyperparameter & Symbol & Value \\
\midrule
Novelty Threshold & $\tau_N$ & 0.50 \\
Cohesion Threshold & $\tau_C$ & 0.15 \\
Confidence Threshold & $\tau_{conf}$ & 0.90 \\
EMA Momentum & $\alpha$ & 0.99 \\
Contrastive Temperature & $\tau$ & 0.10 \\
Learning Rate & $\eta$ & 0.001 \\
Batch Size & $b$ & 64 \\
\bottomrule
\end{tabular}
\end{table}


\end{document}